% !TeX program = xelatex
% Compile this file on its own for a two-page proof.
%
% To use the objects independently in the thesis, put these three lines in
% main.tex before \begin{document} (adjust the input path if needed):
%   \usepackage{tikz}
%   \newcommand{\ThesisLiteratureEmbedded}{}
%   \input{review_figure_table.tex}
% Then place \ReviewDepthFigure and \ReviewFollowupTable at their respective
% locations in Chapter 2. Cross-references: fig:review-depth, tab:review-followup.
% The thesis already loads fontspec, geometry, setspace, array, tabularx,
% booktabs, float, caption and biblatex. Its caption and font settings apply.
%
% The values describe the original appraisal records. Changing the question
% to coverage in the full papers would require recoding before using it here.
\ifdefined\ThesisLiteratureEmbedded\else
  \documentclass[12pt,a4paper,oneside]{report}
  \usepackage{fontspec}
  \usepackage[margin=1in]{geometry}
  \usepackage{setspace,array,tabularx,booktabs,float,caption,tikz}
  \captionsetup{labelfont=bf,textfont=it,labelsep=newline,
    justification=raggedright,singlelinecheck=false,position=top}
  % Use neutral preview numbering. The thesis supplies its own numbers.
  \pagestyle{empty}
\fi

% DATA -----------------------------------------------------------------------
% Arguments: topic ID, label, direct test, used/discussed, no record, unclear.
% One status per topic per appraisal; a work can contribute to several topics.
% The figure and table both read this block.
\newcommand{\ReviewAppraisalN}{94}
\newcommand{\ReviewTopics}[1]{%
  #1{inputs}{Adding an input}{46}{15}{30}{3}%
  #1{rule}{Choice of fusion rule}{37}{28}{29}{0}%
  #1{training}{Training strategy}{28}{60}{4}{2}%
  #1{noise}{Corrupted inputs at test}{21}{22}{50}{1}%
  #1{encoding}{Input representation}{15}{67}{7}{5}%
  #1{absent}{Missing inputs at inference}{11}{15}{68}{0}%
  #1{shift}{Transfer across settings}{11}{30}{52}{1}%
}
\newcommand{\ReviewPercent}[1]{%
  \begingroup
  \pgfmathparse{100*(#1)/\ReviewAppraisalN}%
  \pgfmathprintnumber[fixed,precision=1,fixed zerofill]{\pgfmathresult}\%%
  \endgroup}
\newcommand{\ReviewCheckRow}[6]{%
  \ifnum\numexpr #3+#4+#5+#6\relax=\ReviewAppraisalN\relax\else
    \PackageError{review-visuals}{Counts for #1 must sum to \ReviewAppraisalN}%
      {Check the four status counts in the ReviewTopics block.}%
  \fi}
\ReviewTopics{\ReviewCheckRow}

% Native author-year citations in the thesis; the stand-alone proof prints
% the same work identities without needing references.bib or a Biber run.
\newcommand{\ReviewCite}[2]{%
  \ifcsname parencite\endcsname\parencite{#1}\else(#2)\fi}
\newcommand{\ReviewNote}[1]{%
  \par\smallskip\begingroup\small\raggedright\noindent\textit{Note.} #1\par\endgroup}

% FIGURE ---------------------------------------------------------------------
% Coordinates are fractions of the current text width; row spacing is in cm.
% Greyscale fills and hatching remain distinguishable in monochrome print.
\newcount\ReviewRowIndex
\newcommand{\ReviewBarRow}[6]{%
  \advance\ReviewRowIndex by 1\relax
  \pgfmathsetmacro{\ReviewY}{-0.72*\the\ReviewRowIndex}%
  \pgfmathsetmacro{\ReviewA}{\ReviewLeft+\ReviewWidth*(#3)/\ReviewAppraisalN}%
  \pgfmathsetmacro{\ReviewB}{\ReviewLeft+\ReviewWidth*(#3+#4)/\ReviewAppraisalN}%
  \pgfmathsetmacro{\ReviewC}{\ReviewLeft+\ReviewWidth*(#3+#4+#5)/\ReviewAppraisalN}%
  \node[anchor=west] at (0,\ReviewY) {#2};
  \path[fill=black!76] (\ReviewLeft,{\ReviewY-.18}) rectangle (\ReviewA,{\ReviewY+.18});
  \path[fill=black!36] (\ReviewA,{\ReviewY-.18}) rectangle (\ReviewB,{\ReviewY+.18});
  \path[fill=white] (\ReviewB,{\ReviewY-.18}) rectangle (\ReviewC,{\ReviewY+.18});
  \ifnum#6>0\relax
    \begin{scope}
      \clip (\ReviewC,{\ReviewY-.18}) rectangle (\ReviewRight,{\ReviewY+.18});
      \pgfmathtruncatemacro{\ReviewHatches}{ceil((\ReviewRight-\ReviewC)/.008)+3}%
      \foreach \j in {0,...,\ReviewHatches} {
        \draw[black!75,line width=.28pt]
          ({\ReviewC-.020+.008*\j},{\ReviewY-.18}) --
          ({\ReviewC+.002+.008*\j},{\ReviewY+.18});
      }
    \end{scope}
    \draw[black!55,line width=.25pt] (\ReviewC,{\ReviewY-.18}) -- (\ReviewC,{\ReviewY+.18});
  \fi
  \draw[black!65,line width=.3pt] (\ReviewLeft,{\ReviewY-.18}) rectangle (\ReviewRight,{\ReviewY+.18});
  \node at (.745,\ReviewY) {#3}; \node at (.818,\ReviewY) {#4};
  \node at (.889,\ReviewY) {#5}; \node at (.958,\ReviewY) {#6};
}
\newcommand{\ReviewKey}[1]{%
  \tikz[baseline=-.5ex,x=1pt,y=1pt]{%
    \ifcase#1\relax\or \fill[black!76] (0,0) rectangle (11,6);
    \or \fill[black!36] (0,0) rectangle (11,6);
    \or \fill[white] (0,0) rectangle (11,6);
    \or \begin{scope}\clip (0,0) rectangle (11,6);
      \foreach \j in {-6,-3,0,3,6,9}{\draw[line width=.3pt] (\j,0)--++(6,6);}
    \end{scope}\fi
    \draw[black!65,line width=.3pt] (0,0) rectangle (11,6);}}
\newcommand{\ReviewDepthFigure}{%
  \begin{figure}[H]
    \centering\singlespacing
    \caption{Depth of investigation recorded in the \ReviewAppraisalN{} literature appraisals}
    \label{fig:review-depth}
    \begingroup\footnotesize
    \begin{tabular}{@{}l@{\hspace{2em}}l@{}}
      \ReviewKey{1}\enspace Direct test recorded &
      \ReviewKey{2}\enspace Used or discussed only\\[2pt]
      \ReviewKey{3}\enspace No treatment recorded &
      \ReviewKey{4}\enspace Unclear from appraisal
    \end{tabular}\par\medskip
    \begin{tikzpicture}[x=\linewidth,y=1cm,
      every node/.style={inner sep=0pt,outer sep=0pt,font=\footnotesize}]
      \def\ReviewLeft{.300}\def\ReviewWidth{.388}\def\ReviewRight{.688}
      \node[anchor=west,font=\footnotesize\bfseries] at (0,.53) {Topic};
      \node[font=\footnotesize\bfseries] at (.494,.53) {Share of appraisals};
      \node[font=\footnotesize\bfseries] at (.854,.53) {Number of appraisals};
      \foreach \tick in {0,25,50,75,100}{
        \pgfmathsetmacro{\xx}{\ReviewLeft+\ReviewWidth*\tick/100}
        \node[font=\scriptsize] at (\xx,.05) {\tick\%};
        \draw[black!20,line width=.25pt] (\xx,-.20) -- (\xx,-5.36);
      }
      \node at (.745,.05) {Test}; \node at (.818,.05) {Use};
      \node at (.889,.05) {None}; \node at (.958,.05) {Unclear};
      \ReviewRowIndex=0\relax
      \ReviewTopics{\ReviewBarRow}
    \end{tikzpicture}
    \endgroup
    \ReviewNote{Each bar counts all \ReviewAppraisalN{} appraisals. A direct test requires
      a recorded comparison or intervention. The same work may address several topics.
      The figure is based on the author's literature appraisals. OpenAI Codex assisted with organising
      the records, calculating the percentages and preparing the figure.}
  \end{figure}}

% TABLE ----------------------------------------------------------------------
% Each topic has two editable cells: recorded finding and author-proposed
% follow-up. Citation keys in the first argument of ReviewCite are those used
% in the thesis bibliography. The second argument is the stand-alone fallback.
\newcommand{\ReviewText}[3]{%
  \expandafter\def\csname ReviewFinding#1\endcsname{#2}%
  \expandafter\def\csname ReviewQuestion#1\endcsname{#3}}
\ReviewText{inputs}
  {Image and audio fusion gained 5.9 percentage points over image alone in four-state welding classification.
   In weld penetration, each single input tied fusion on the clean test
   \ReviewCite{fan_dmaw_2026,hua_anti-interference_2026}{Fan et al., 2026; Hua et al., 2026}.}
  {Does a second input improve on the strongest single input when both are tested on the same held-out runs?}
\ReviewText{rule}
  {An RGB--3D inspection study compared a static mean with a learned gate using complete inputs
   \ReviewCite{deng_tc-maf_2026}{Deng et al., 2026}.}
  {Under which faults does a learned rule outperform a static rule after both receive the same exposure during training?}
\ReviewText{training}
  {Both models had 1.02\% error on clean digits. With two quadrants missing, modality-dropout training reduced error from 35.91\% to 7.19\%
   \ReviewCite{neverova_moddrop_2015}{Neverova et al., 2015}.}
  {How much of a gain under corruption comes from training exposure, and how much from the fusion rule?}
\ReviewText{noise}
  {One bearing study compared fusion rules at 0~dB. Another reported the proposed model's accuracy when one input was degraded
   \ReviewCite{huang_rolling_2026,P046}{Huang \& Tong, 2026; Sun et al., 2026}.}
  {At the same fault severity, does fusion outperform both the intact input and a static combination?}
\ReviewText{encoding}
  {Handcrafted geometry and learned video features were evaluated on different targets
   \ReviewCite{horwitz_back_2023,liu_video_2021}{Horwitz \& Hoshen, 2023; Liu et al., 2021}.}
  {Which encoding adds complementary information when the target, observation window and fitting budget are fixed?}
\ReviewText{absent}
  {The appraisals distinguish planned single-input operation, simulated dropout and missing 3D data
   \ReviewCite{li_pgmtkd_2025,xu_modality-resilient_2025,deng_tc-maf_2026}{Li et al., 2025; Xu et al., 2025; Deng et al., 2026}.}
  {Does a policy trained for a known missing stream also cope with an unexpected outage during operation?}
\ReviewText{shift}
  {Handover was tested across robots. A bearing study used target-domain validation during transfer
   \ReviewCite{thoduka_multimodal_2024,huang_rolling_2026}{Thoduka et al., 2024; Huang \& Tong, 2026}.}
  {How much fusion benefit persists on a new machine or site when access to target data is controlled?}
\newcommand{\ReviewTableRow}[6]{%
  \textbf{#2}\newline #3/\ReviewAppraisalN{} (\ReviewPercent{#3}) &
  \csname ReviewFinding#1\endcsname & \csname ReviewQuestion#1\endcsname\tabularnewline
  \addlinespace[.45em]}
\newcommand{\ReviewFollowupTable}{%
  \begin{table}[H]
    \centering\singlespacing
    \caption{Recorded findings and follow-up questions from the literature review}
    \label{tab:review-followup}
    \begingroup\small
    \setlength{\tabcolsep}{4.5pt}
    \renewcommand{\arraystretch}{1.04}
    \begin{tabularx}{\linewidth}{@{}>{\raggedright\arraybackslash}p{.207\linewidth}
      >{\raggedright\arraybackslash}X >{\raggedright\arraybackslash}X@{}}
      \toprule
      \textbf{Topic and coverage} & \textbf{Evidence recorded in the appraisal}
        & \textbf{Follow-up questions suggested by the review}\\
      \midrule
      \ReviewTopics{\ReviewTableRow}
      \bottomrule
    \end{tabularx}
    \endgroup
    \ReviewNote{Coverage gives the number and share of the \ReviewAppraisalN{} appraisals
      recording a direct test. The last column gives the author's interpretation of the
      cited comparisons. OpenAI Codex assisted with preparing the table.}
  \end{table}}

% Stand-alone preview only. In the thesis, each object is placed by its macro.
\ifdefined\ThesisLiteratureEmbedded
  \expandafter\endinput
\fi
\begin{document}
\ReviewDepthFigure
\clearpage
\ReviewFollowupTable
\end{document}
